\chapter{Introducción}

La realidad virtual (VR, por sus siglas en inglés) se ha 
consolidado como un campo de investigación relevante dentro del 
desarrollo de experiencias digitales inmersivas, sobre todo en lo
relacionado con el diseño de interfaces, técnicas de interacción
y evaluación de la experiencia de 
usuario \cite{chen2024interfaces}. Estudios recientes muestran 
que la interacción en entornos virtuales no puede abordarse 
únicamente desde una perspectiva técnica, sino que requiere 
considerar criterios de usabilidad, naturalidad, 
retroalimentación y un contexto adecuado para la experiencia 
\cite{chen2024interfaces,spittle2023interaction,yeo2024interfaces}. Con eso 
en mente, este proyecto propone la implementación de un sistema de 
interacción adaptativa en VR para una experiencia de videoclip gamificado. 
La propuesta se enfoca en el desarrollo de una solución capaz de registrar 
métricas de desempeño del usuario en tiempo real, interpretar estados de 
interacción y ajustar dinámicamente parámetros del \textit{gameplay}, con el
fin de construir una experiencia más sensible al comportamiento del jugador 
y una mejor integridad al flujo general de la aplicación.

El valor de este proyecto se sustenta en investigaciones recientes sobre 
adaptación en videojuegos y entornos interactivos han señalado que, aunque 
existen avances importantes en personalización, modelado del jugador y 
ajuste dinámico del contenido, aún queda una brecha en la integración 
completa del ciclo conformado por las sensaciones del usuario, la 
interpretación de su estado y la modificación adaptativa de la experiencia 
en tiempo real \cite{lopes2025closing}. Debido a esto, resulta académica y
técnicamente relevante desarrollar prototipos que articulen la telemetría 
del usuario con mecanismos de adaptación concretos dentro de experiencias 
inmersivas. A su vez, investigaciones sobre ajustes dinámicos de dificultad 
en VR evidencian que la adaptación puede ir más allá de simples cambios de 
nivel, permitiendo modificar velocidad, complejidad y estructura de las 
tareas de acuerdo con las capacidades o desempeño del usuario 
\cite{huber2021dda}. Esto resulta particularmente valioso en una experiencia 
de videoclip gamificado, donde el equilibrio entre reto, ritmo, claridad e
involucramiento constituye un factor central para sostener la atención del 
jugador y favorecer una participación más fluida.

De igual manera, dentro de la industria, la relevancia que adquiere el 
proyecto se relaciona con la creciente necesidad de diseñar experiencias 
inmersivas que no sean estáticas, sino capaces de responder al usuario de 
manera contextual \cite{lopes2025closing,huber2021dda}. La revisión 
sistemática de Lopes et al. señala que los sistemas adaptativos orientados 
por la experiencia del jugador representan una dirección prometedora para 
incrementar la personalización y la inmersión en aplicaciones interactivas
\cite{lopes2025closing}. Las investigaciones sobre evaluación en VR también 
reportan que constructos como presencia, experiencia de usuario y mareos 
virtuales constituyen dimensiones centrales para valorar la calidad de una 
experiencia inmersiva, y que su análisis suele beneficiarse de la 
combinación entre métricas objetivas y herramientas subjetivas de 
evaluación \cite{wang2025presence,bareisyte2024questionnaires}. De la misma 
manera, se ha documentado que el \textit{feedback} visual y háptico puede contribuir 
positivamente a la sensación de presencia, mientras que la carga de tareas 
y las exigencias impuestas al usuario pueden influir en el nivel de malestar 
o mareos virtuales \cite{sepich2022workload,gibbs2022feedback}. Por ello, 
el proyecto no solo resulta importante por su componente técnico, sino 
también por su potencial para explorar cómo la adaptación dinámica puede 
incidir en la calidad de la experiencia inmersiva.

Como antecedentes, se han desarrollado diferentes propuestas que demuestran 
la viabilidad de utilizar información observable del jugador para modificar 
una experiencia interactiva durante su ejecución. Algunos sistemas han 
construido modelos a partir del comportamiento registrado durante la partida
para seleccionar modificaciones orientadas a mejorar la experiencia estimada
del usuario \cite{yannakakis2009adaptation}. Otros trabajos han considerado 
la dificultad como un fenómeno temporal, analizando tendencias a lo largo de
varias interacciones antes de decidir si el desafío debe incrementarse, 
mantenerse o reducirse \cite{zook2012temporal,xue2017dda}. Dentro del contexto específico de VR, también se han implementado mecanismos de ajuste
dinámico capaces de modificar la estructura y complejidad de las tareas de acuerdo con
información obtenida del usuario \cite{huber2021dda}. Estas aproximaciones 
evidencian que el desempeño del jugador puede utilizarse como una variable 
dinámica y no únicamente como un resultado evaluado al finalizar una sesión. 
También muestran que la adaptación en VR puede intervenir directamente sobre 
variables del \textit{gameplay} y no limitarse a niveles generales de 
dificultad. Paralelamente, los videojuegos musicales han desarrollado 
distintas formas de integrar la música con las acciones del jugador, 
utilizando elementos sonoros y temporales como parte de las propias reglas 
de interacción \cite{pichlmair2007sound}. Sin embargo, la integración 
conjunta de una experiencia musical en VR con un sistema que registre y 
modifique parámetros del \textit{gameplay} en tiempo real representa una 
convergencia menos desarrollada.

El proyecto se desarrolla en un lapso de cinco meses y está orientado a la 
construcción de un prototipo académico funcional en Unity, utilizando 
Blender para la creación de algunos recursos tridimensionales. El desarrollo
se enfoca en la interacción base en VR, el registro de métricas de desempeño, 
el modelado de estados de interacción y la adaptación dinámica del gameplay.
Para ello, se consideran las siguientes métricas: errores, tiempo de 
reacción, precisión y frecuencia de interacción. Estas métricas permiten 
observar tendencias en el comportamiento del jugador durante la ejecución de
la experiencia. A partir de ellas, el sistema utilizará estados discretos como 
representaciones computacionales del comportamiento observado y podrá 
modificar parámetros directamente relacionados con el \textit{gameplay}, 
principalmente velocidad, cantidad de objetos y complejidad de las tareas. 
Estos estados no pretenden construir diagnósticos psicológicos o clínicos, 
sino abstracciones operativas orientadas a la toma de decisiones dentro del 
sistema adaptativo.

El alcance contempla también los mecanismos necesarios para producir y 
registrar dichas interacciones, entre ellos entrada mediante controles, 
detección de colisiones, reglas de juego y retroalimentación visual y háptica. 
Por otra parte, quedan fuera del desarrollo los componentes relacionados con
la secuenciación temporal global del videoclip, la administración general
del tiempo de la canción, la activación global de escenas, la sincronización
audiovisual general y la persistencia del progreso. Estos elementos podrán 
comunicarse con el sistema adaptativo, pero no forman parte del mecanismo de 
medición, modelado y ajuste desarrollado en este trabajo. Asimismo, se 
excluyen la comercialización del producto final y cualquier validación 
clínica de estados cognitivos complejos.

Las pruebas realizadas con 34 participantes generan 5048 registros de 
interacción distribuidos entre los seis minijuegos de la experiencia y 
muestran que el sistema es capaz de registrar diferencias de desempeño, 
representar cambios mediante estados operativos y modificar la 
dificultad durante la ejecución. El sistema de adapatación dinámica de 
dificultad (DDA, por sus siglas en inglés) realiza 216 cambios de 
dificultad y, entre los ajustes que permiten comparar cinco 
interacciones antes y después, las reducciones presentan un aumento 
promedio en la tasa de éxito de 34.71\% a 62.06\%, mientras que los 
aumentos de dificultad producen una reducción de 91.30\% a 60.00\%, 
evidenciando un cambio observable en la exigencia de las actividades. 
En la evaluación subjetiva, el 67.65\% de los participantes reporta 
niveles altos de atención a la canción, el 96.97\% de las respuestas 
válidas califica la diversión con valores de 4 o 5 y el 73.53\% indica 
sentirse nada o ligeramente abrumado. Estos resultados respaldan el 
funcionamiento del sistema de interacción adaptativa, aunque no permiten 
atribuir directamente las valoraciones subjetivas al DDA.

